\section{Estimation and inference choices}
\label{sec:modern}

An estimator identifies nothing on its own; the design does. The choices in this section can protect or squander what the design buys, but they can't rescue a contrast whose counterfactual is implausible.

The default estimator is a trap. A pooled two-way fixed-effects regression looks natural, but under staggered adoption with heterogeneous effects its implied weights can turn negative, so already-treated units serve as controls for later ones and can push the overall estimate the wrong way \citep{GoodmanBacon2021,deChaisemartin2020}. Group-time and event-study estimators repair that weighting problem \citep{CallawaySantAnna2021,SunAbraham2021,BorusyakJaravelSpiess2024}. They don't repair selection into timing: a heterogeneity-robust estimate of an endogenously timed treatment is still a clean number for the wrong quantity.

Scale is substantive too. The additive split $f = \pi + b$ buys clean linearity, but parallel trends isn't transform-invariant: an assumption that holds on the probability scale can fail on the log-odds scale, especially when baseline rates differ \citep{RothSantAnna2023}. A log-odds model may be the more defensible estimator for cell rates, and beta-binomial models may handle overdispersion, but the effect should be reported on the scale on which the claim is made.

Inference is where the few-units problem bites. The binding sample size is the number of identifying units, not the millions of tokens, and with four or five polities ordinary cluster-robust intervals are badly anti-conservative \citep{BertrandDufloMullainathan2004}. The usual repairs are the wild cluster bootstrap \citep{CameronGelbachMiller2008}, resampling whole varieties, and randomization inference for few treated clusters \citep{MacKinnonWebb2020}.

Neither repair is a guarantee: the bootstrap can still under-cover with very few or unbalanced clusters, and randomization inference needs a credible assignment model that politically chosen adoption rarely supplies. Where their assumptions can't be made explicit, the absence of a clean interval is part of the result.

The unit of analysis should be pre-specified cells, such as outlet-period or title-strategy cells, not raw tokens treated as independent draws. This choice is the aggregation choice of Section~\ref{sec:estimand} seen from the inference side.

This raises a modelling objection: why difference at all, rather than model the whole series with latent rates, outlet and title effects, composition terms, and a smooth calendar trend? Differencing clears time-invariant confounds but leaves trend confounding exactly where a model would: selection into treatment timing rides the trend, and no estimator removes it. DiD's advantage is a named and conventional assumption. Parallel trends sits in plain view, and a flat pre-period is its agreed probe; a multilevel model can bury the same assumption in functional form and random-effects structure.

Partial pooling is still useful for sparse outlet and title cells within each variety; it's a complement for estimation, not a substitute for the design's identification. None of these choices identifies anything. They belong after the estimand and diagnostics, and before the worked example puts them to work.
